\subsection{层的子层}

设 B 为一个拓扑空间。设 \(\mathcal{F} = (\mathcal{F}(U), f_{UV})\) 为 B 上的、相对于 B 的拓扑的一个基 \(\mathcal{B}\) 的预层。

假设对每个开集 \(U \in \mathcal{B}\)，存在一个子集 \(\mathcal{L}(U)\) （\(\mathcal{F}(U)\) 的）。若 \(f_{\mathrm{UV}}(\mathcal{L}(V)) \subset \mathcal{L}(U)\) 对每对元素 \((U, V)\) （\(\mathcal{B}\) 中的，满足 \(U \subset V\) 者）成立，则对 \(\mathcal{L} = ((\mathcal{L}(U))_{U \in \mathcal{B}}, (f'_{UV}))\) （其中 \(f'_{UV} \colon \mathcal{L}(V) \to \mathcal{L}(U)\) 是由 \(f_{UV}\) 诱导的映射）是一个预层。这样的预层称为 \(\mathcal{F}\) 的子预层。由于诸映射 \(f'_{UV}\) 由预层 F 的数据和族 \((\mathcal{L}(U))_{U \in \mathcal{B}}\) 所确定，人们也滥用语言地说族 \((\mathcal{L}(U))_{U \in \mathcal{B}}\) 是 F 的一个子预层。

现在假设 \(\mathcal{F}\) 是 B 上的一个层，并且对 B 的每个开子集 U，设 \(\mathcal{L}(\mathrm{U})\) 为 \(\mathcal{F}(\mathrm{U})\) 的一个子集。要使 \((\mathcal{L}(\mathrm{U}))_{\mathrm{U}\in \mathcal{B}}\) 是 \(\mathcal{F}\) 的子预层，且要使该预层是一个层，必须且只需下列条件被满足：

\begin{enumerate}
\def\labelenumi{(\Alph{enumi})}
\setcounter{enumi}{5}
\tightlist
\item
  设 \((\mathrm{U}_{i})_{i\in\mathrm{I}}\) 为 B 的开集的一个族，设 U 为其并，并设 s 为 \(\mathcal{F}(\mathrm{U})\) 的一个元素。要使 s 属于 \(\mathcal{L}(\mathrm{U})\) ，必须且只需：
\end{enumerate}

\[
\text{对每个 } i \in \mathrm{I},\quad f_{\mathrm{U}_{i}\mathrm{U}}(s)\in\mathcal{L}(\mathrm{U}_{i}).
\]

事实上，若条件 (F) 成立，则对 B 的满足 U ⊂ V 的每对开集 (U, V)，有 \(f_{\mathrm{UV}}(\mathcal{L}(\mathrm{V})) \subset \mathcal{L}(\mathrm{U})\) ，而性质 \(F_{1}\) 与 \(F_{2}\) （相对于子预层 \(\mathcal{L}(\mathrm{U})\) 者）可由相对于层 F 的类似性质得出。逆命题是显然的。

当条件 (F) 满足时，称 \((\mathcal{L}(\mathrm{U}))\) 为层 F 的子层。

